% !TEX root = ../Main.tex
\chapter{概率（三）：独立重复试验与事件独立}

\ex[例 3]{$A$、$B$ 是治疗同一种疾病的两种药，用若干试验组做对比实验。每个试验组由 $4$ 只小白鼠组成，其中 $2$ 只服用 $A$、另外 $2$ 只服用 $B$。一段时间后观察疗效：若在一个试验组中，服用 $A$ 有效的只数\textbf{比}服用 $B$ 有效的只数\textbf{多}，就称该试验组为\textbf{甲类组}。已知每只小白鼠服用 $A$ 有效的概率为 $\dfrac{2}{3}$，服用 $B$ 有效的概率为 $\dfrac{1}{2}$。

\begin{enumerate}
    \item 求一个试验组为甲类组的概率；
    \item 观察 $3$ 个试验组，求这 $3$ 个组中\textbf{至少有一个}甲类组的概率。
\end{enumerate}}

\sol{
    设 $X$ 为该组中服 $A$ 有效的只数，$Y$ 为服 $B$ 有效的只数。$X$、$Y$ 各是两次独立试验里有效的次数：
    \[
    P(X=k)=C_2^k\pr{\tfrac{2}{3}}^k\pr{\tfrac{1}{3}}^{2-k},\qquad
    P(Y=k)=C_2^k\pr{\tfrac{1}{2}}^2.
    \]
    列出分布：
    \[
    \ba
    &P(X=0)=\tfrac{1}{9},\quad P(X=1)=\tfrac{4}{9},\quad P(X=2)=\tfrac{4}{9};\\
    &P(Y=0)=\tfrac{1}{4},\quad P(Y=1)=\tfrac{1}{2},\quad P(Y=2)=\tfrac{1}{4}.
    \ea
    \]

    \textbf{(1)} 甲类组就是 $X>Y$，把所有 $X>Y$ 的情形列出来（$X$、$Y$ 相互独立，概率相乘）：
    \[
    \ba
    P(X>Y)
    &=P(X{=}1)P(Y{=}0)+P(X{=}2)P(Y{=}0)+P(X{=}2)P(Y{=}1)\\
    &=\frac{4}{9}\cdot\frac{1}{4}+\frac{4}{9}\cdot\frac{1}{4}+\frac{4}{9}\cdot\frac{1}{2}
    =\frac{4}{36}+\frac{4}{36}+\frac{8}{36}=\frac{16}{36}=\frac{4}{9}.
    \ea
    \]
    所以一个试验组为甲类组的概率是 $\dfrac{4}{9}$。

    \textbf{(2)} $3$ 个组相互独立，"至少有一个甲类组"的对立事件是"一个甲类组都没有"：
    \[
    P(\text{至少一个})=1-\pr{1-\frac{4}{9}}^3=1-\pr{\frac{5}{9}}^3
    =1-\frac{125}{729}=\frac{604}{729}.
    \]
}

\ex[例 4]{有 $6$ 个相同的球，分别标有数字 $1,2,3,4,5,6$。从中\textbf{有放回}地随机取两次，每次取 $1$ 个球。记：甲 $=$"第一次取出的球的数字是 $1$"，乙 $=$"第二次取出的球的数字是 $2$"，丙 $=$"两次取出的数字之和是 $8$"，丁 $=$"两次取出的数字之和是 $7$"。则（\quad）

\quad A.\ 甲、丙相互独立\qquad B.\ 甲、丁相互独立

\quad C.\ 乙、丙相互独立\qquad D.\ 丙、丁相互独立}

\sol{
    有放回取两次、记顺序，基本事件共 $6\times 6=36$ 个，等可能。先各算概率：
    \[
    P(\text{甲})=\frac{6}{36}=\frac{1}{6},\qquad
    P(\text{乙})=\frac{6}{36}=\frac{1}{6}.
    \]
    和为 $8$ 的有 $(2,6),(3,5),(4,4),(5,3),(6,2)$ 共 $5$ 种，和为 $7$ 的有 $(1,6),(2,5),(3,4),(4,3),(5,2),(6,1)$ 共 $6$ 种：
    \[
    P(\text{丙})=\frac{5}{36},\qquad P(\text{丁})=\frac{6}{36}=\frac{1}{6}.
    \]
    判独立就看 $P(\text{交})$ 是否等于两概率之积。

    \begin{itemize}
        \item 甲丙：第一次为 $1$、和为 $8$ 需第二次为 $7$，不可能，$P(\text{甲}\cap\text{丙})=0\ne\dfrac{1}{6}\cdot\dfrac{5}{36}$，\textbf{不独立}；
        \item 甲丁：只有 $(1,6)$，$P(\text{甲}\cap\text{丁})=\dfrac{1}{36}=\dfrac{1}{6}\cdot\dfrac{1}{6}=P(\text{甲})P(\text{丁})$，\textbf{独立}；
        \item 乙丙：只有 $(6,2)$，$P(\text{乙}\cap\text{丙})=\dfrac{1}{36}$，而 $P(\text{乙})P(\text{丙})=\dfrac{1}{6}\cdot\dfrac{5}{36}=\dfrac{5}{216}\ne\dfrac{6}{216}$，\textbf{不独立}；
        \item 丙丁：和不能同时是 $8$ 又是 $7$，$P(\text{丙}\cap\text{丁})=0\ne P(\text{丙})P(\text{丁})$，\textbf{不独立}（它俩其实互斥）。
    \end{itemize}

    只有 B 满足 $P(AB)=P(A)P(B)$，选 \textbf{B}。
}
